\documentclass[11pt]{article}
\usepackage[margin=25mm]{geometry}
\usepackage{amsmath,amssymb,amsthm}
\newtheorem{theorem}{Theorem}
\newtheorem{lemma}{Lemma}
\newtheorem{corollary}{Corollary}
\title{A bound between singleton distance layers}
\author{Human Compression project}
\date{}
\begin{document}\maketitle

Let $G$ be a finite connected simple four-regular graph of girth at least
five. Fix a vertex and let $V_0,\ldots,V_T$ be its distance layers.
Put $n_h=|V_h|$. All distances and layers in this note refer to this fixed
root. The argument uses no assumption about the diameter of $G$.

\begin{theorem}\label{thm:section}
Suppose that $1\leq i$, $i+5<T$, and $n_i=n_{i+5}=1$. Then
\[
 n_{i+1}=2,\qquad n_{i+2}=6,\qquad
 n_{i+1}+n_{i+2}+n_{i+3}+n_{i+4}\geq18.
\]
\end{theorem}

We first record the elementary restrictions that will be used in the proof.
The edge boundary of any vertex set is even, since
$|\partial S|=4|S|-2e(G[S])$. Every proper nonempty distance prefix has a
nonempty boundary, by connectedness. Consequently a singleton layer
$V_h=\{w\}$ with $0<h<T$ has exactly two edges to each of its neighboring
layers. There are no edges between layers whose indices differ by more
than one.

\begin{lemma}\label{lem:parents}
Suppose that an internal singleton layer $V_h$, with $0<h<T$, is followed
by layers $A,B,C$, in that order. The first two have sizes $2,6$.
The vertices of $B$ split into two
sets $P=\{p_0,p_1,p_2\}$ and $Q=\{q_0,q_1,q_2\}$ according to their
parent in $A$. The edges within $B$ form a matching between $P$ and $Q$.
Each vertex of $C$ has at most one neighbor in each of $P,Q$.

If this matching has size three, no two vertices of $C$ which each have
two neighbors in $B$ can be adjacent.
\end{lemma}
\begin{proof}
The singleton has two distinct neighbors in $A$. They cannot be adjacent,
and their other three neighbors lie in $B$. These six neighbors are
distinct: a shared neighbor would give a quadrangle through the singleton.
Every vertex of $B$ has a parent in $A$, so these are all of $B$.

An edge within either parent class would give a triangle. If a vertex
had two neighbors in the other parent class, those two neighbors and
their common parent would give a quadrangle. Thus the internal edges
form the claimed matching. Two neighbors of a vertex of $C$ in the
same parent class likewise give a quadrangle with their common parent.

For the last assertion, number the vertices so that the matching is
$p_jq_j$ for $j=0,1,2$. A vertex of $C$ with two parents represents a
pair $(a,b)$, meaning that its neighbors in $B$ are $p_a,q_b$.
We have $a\ne b$, since otherwise there is a triangle. If two such
vertices, representing $(a,b)$ and $(c,d)$, were adjacent, the absence
of triangles would imply $a\ne c$ and $b\ne d$. The absence of
quadrangles using an edge of the matching would also imply
$a\ne d$ and $b\ne c$. Both $c$ and $d$ would therefore be the unique
index outside $\{a,b\}$, contrary to $c\ne d$.
\end{proof}

\begin{proof}[Proof of Theorem~\ref{thm:section}]
Write $V_i=\{u\}$, $V_{i+5}=\{v\}$, and name the intervening layers
$A,B,C,D$. By Lemma~\ref{lem:parents}, $|A|=2$, $|B|=6$. Put
$c=|C|$, $d=|D|$, and let $m$ be the size of the matching inside $B$.
The number of edges from $B$ to $C$ is
\[
 e(B,C)=4\cdot6-6-2m=18-2m.
\]
Each vertex of $C$ has at most two parents in $B$, and $m\leq3$.
It follows that $c\geq6$. The internal singleton $v$ has exactly two
edges from $D$, so $d\geq2$.

Suppose, for a contradiction, that the four intervening layers have
at most $17$ vertices. Then $c+d\leq9$, leaving only
\[
 (c,d)=(6,2),\ (6,3),\ (7,2).
\]
If $c=6$, the inequality $18-2m\leq12$ forces $m=3$ and every vertex
of $C$ has two parents. Lemma~\ref{lem:parents} says that $C$ is
independent. Its remaining twelve incidences must go to $D$.
But $d\leq3$ and two incidences of $D$ go to $v$, giving
$e(C,D)\leq4d-2\leq10$, a contradiction.

It remains to consider $(c,d)=(7,2)$. Both vertices of $D$ are adjacent
to $v$, so they are nonadjacent. Each has three neighbors in $C$,
and the two triples are disjoint, since a shared neighbor would give
a quadrangle through $v$. Denote these triples by $C_1,C_2$, and let
$z$ be the seventh vertex of $C$. Thus $e(C,D)=6$.
Degree counting in $C$ gives
\[
 28=(18-2m)+2e(G[C])+6,
 \qquad e(G[C])=m+2.
\]
Since $18-2m\leq14$, we have $m=2$ or $m=3$.

If $m=3$, then $e(G[C])=5$. There are no edges within $C_1$ or $C_2$.
Every vertex has at most one neighbor in each of these triples: two
would give a quadrangle through the corresponding vertex of $D$.
Hence the edges between $C_1,C_2$ form a matching of size at most three,
and $z$ has at most one neighbor in each triple. Equality in the bound
$e(G[C])\leq3+2=5$ forces a matching of size three and two edges at $z$.
The two neighbors of $z$ cannot be matched to each other, since that
would give a triangle. They therefore meet two distinct matching edges.
Their degrees within $C$ are two, the other four vertices in the triples
have internal degree one, and $z$ has internal degree two.

Each vertex of the triples has one neighbor in $D$. Consequently the
two neighbors of $z$ have one parent in $B$, while the other four
vertices and $z$ have two. Of the three matching edges between the
triples, at least one avoids the two vertices with only one parent.
Its endpoints both have two parents in $B$, contradicting
Lemma~\ref{lem:parents}.

Finally suppose $m=2$. Number the matching edges in $B$ as
$p_0q_0,p_1q_1$, leaving $p_2,q_2$ unmatched. There are fourteen edges
from $B$ to the seven vertices of $C$, so each has two parents.
The seven pairs identifying these vertices are precisely
\[
 (0,1),(1,0),(0,2),(2,0),(1,2),(2,1),(2,2).
\]
Indeed $(0,0)$ and $(1,1)$ would give triangles, and a repeated pair
would give a quadrangle. For two pair labels $(a,b),(c,d)$ to be
adjacent in $C$, they must have different first coordinates and different
second coordinates, and neither $(a,d)$ nor $(c,b)$ may be $(0,0)$
or $(1,1)$. These conditions exclude triangles with a parent and
quadrangles through a matching edge in $B$.
Checking the seven displayed pairs leaves exactly the following possible
edges:
\[
 \{(2,2),(0,1)\},\quad \{(2,2),(1,0)\},\quad
 \{(0,2),(2,1)\},\quad \{(1,2),(2,0)\}.
\]
Since $e(G[C])=m+2=4$, all four must occur. The vertex $(2,2)$ has
internal degree two and hence no neighbor in $D$; the other six each
have internal degree one and one neighbor in $D$.

The three neighbors in $C$ of a fixed vertex of $D$ cannot share any
parent in $B$, since this would give a quadrangle. Their pair labels
must therefore use each first coordinate and each second coordinate
once. Among the six off-diagonal pairs the only such triples are
\[
 \{(0,1),(1,2),(2,0)\},\qquad
 \{(1,0),(0,2),(2,1)\}.
\]
But the first triple contains the internal edge
$\{(1,2),(2,0)\}$ and the second contains
$\{(0,2),(2,1)\}$. Either triple would give a triangle with its
common neighbor in $D$. This is the final contradiction.
\end{proof}

\begin{corollary}
The scalar profile of total order $93$ in
\texttt{layer\_profile\_target\_result/result.json} is not realizable
as the distance layers of the stipulated graph.
\end{corollary}
\begin{proof}
It has $n_9=n_{14}=1$ and
$(n_{10},n_{11},n_{12},n_{13})=(2,6,7,2)$, whose sum is $17$.
Its last layer has index $22$, so Theorem~\ref{thm:section} applies.
\end{proof}

\paragraph{Relation to the existing argument.}
The parity of cuts and the elementary triangle and quadrangle restrictions
are standard and are used here without a novelty claim. The contribution
of this note is their joint application to a section bounded by two
singleton layers, including the two possible matchings in the middle.
The proof uses no computation, external classification, or new group result.
It supplements the existing aggregate inequalities; it does not invalidate
the independently checked feasibility of that weaker scalar model.

\paragraph{Limits.}
This excludes one returned scalar profile and establishes a necessary
condition for every graph meeting the hypotheses. It does not exclude all
graphs of order at most $94$, complete F009, or change the accepted
enumeration prefix. The condition is not yet integrated into the frozen
enumerator or scalar program. No Lean verification is claimed.
\end{document}
